\subsection{جداسازی عروق و پنج نشانگر نهایی}

مسیر نشانگری تصویر را مستقیماً به طبقه‌بند نمی‌دهد. برای هر تصویر یک نقشه‌ی احتمال عروقی ساخته می‌شود، آستانه‌گذاری از آن یک ماسک دودویی می‌سازد و اندازه‌گیری‌ها از همان ماسک استخراج می‌شوند؛ پس معماری جداساز و آستانه‌ی ماسک بخشی از تعریف داده‌ی ورودی این مسیرند.

\subsubsection{جداساز عملیاتی}

جداساز عملیاتی این پروژه شبکه‌ی \lr{MAnet}\enfoot{نام معماری جداساز عروق؛ معادل انگلیسی: \lr{MAnet}} با رمزگذار \lr{ResNet34}\enfoot{شبکه‌ی عصبی باقیمانده با ۳۴ لایه؛ معادل انگلیسی: \lr{ResNet34}} است. تصویر فوندوس به این شبکه داده می‌شود و خروجی احتمال آن با آستانه‌گذاری به ماسک دودویی تبدیل می‌شود. ماسک‌های مصرفی مسیر عروقی از مجموعه‌ی تثبیت‌شده‌ی پروژه می‌آیند تا بازتولید نتایج به یک نسخه‌ی مشخص گره بخورد. بازیابی ماسک فقط نام دقیق ثبت‌شده را می‌پذیرد و هیچ بازگشت تقریبی در آن راه ندارد.

\subsubsection{پنج نشانگر نهایی}

مسیر نشانگری دقیقاً پنج عدد اسکالر تولید می‌کند. \lr{vessel\_density\_fov} نسبت پیکسل‌های عروقی به مساحت محتوا در میدان دید است و \lr{skel\_density\_fov} همان نسبت را برای پیکسل‌های اسکلت (خط مرکزی) عروق می‌سنجد. سه نشانگر دیگر، \lr{fractal\_d0}، \lr{fractal\_d1} و \lr{fractal\_d2}، ابعاد فراکتالی مرتبه‌ی صفر، یک و دو هستند که از تحلیل چندفراکتالی ماسک عروق به‌دست می‌آیند. هر پنج نشانگر روی میدان دید تعریف شده‌اند، نه روی ناحیه‌ای بسته به دیسک بینایی؛ همین تعریف مقایسه‌ی مستقیم آن‌ها را با سامانه‌هایی که ناحیه‌ی بررسی را با قطر دیسک می‌سنجند محدود می‌کند. تعریف اندازه‌گیری در نسخه‌ی تثبیت‌شده‌ی پروژه مخرج نشانگرهای پوشش را برابر مساحت محتوا در میدان دید می‌گذارد و همین تعریف مقیاس اندازه‌گیری را از حاشیه‌های بی‌محتوا جدا می‌کند.

\begin{figure}[htbp]
  \centering
  \imgbox[1.0\textwidth]{figures/report/methods/M2_biomarker_pipeline.pdf}
  \caption{خط لوله‌ی ساخت پنج نشانگر اسکالر از تصویر فوندوس و نقشه‌ی عروق.}
  \label{fig:U4}
\end{figure}


\subsubsection{محدودیت اعتبارسنجی جداسازی}

ماسک مرجع متخصص برای این جمعیت در دسترس نیست؛ بنابراین هیچ معیار هم‌پوشانی نوع \lr{Dice}\enfoot{معیار هم‌پوشانی دو ماسک دودویی؛ معادل انگلیسی: \lr{Dice}} روی داده‌ی خود این پروژه محاسبه نشده است و درباره‌ی صحت بالینی جداسازی در این جمعیت ادعایی مطرح نمی‌شود.

این شکاف تفسیر نتیجه‌ی نشانگرهای اسکالر را محدود می‌کند و دو علت ممکن را از هم جدا نمی‌کند. یک احتمال آن است که پنج عدد اسکالر واقعاً اطلاعات کمتری از بازنمایی یادگرفته‌شده‌ی شبکه حمل می‌کنند؛ احتمال دیگر آن است که کیفیت ماسک و اندازه‌گیری برای این جمعیت محدود باشد و بخشی از ضعف نتیجه از همین‌جا بیاید. طرح آزمایش‌های این گزارش توان جدا کردن این دو علت را ندارد و قضاوت میان آن‌ها به ارزیابی جداسازی با ماسک مرجع نیاز دارد. آنچه داده نشان می‌دهد این است که در رژیم‌های آزمون‌شده، افزودن این پنج نشانگر افزایش پشتیبانی‌شده‌ای فراتر از بازنمایی یادگرفته‌شده ایجاد نکرد.
